\section{Flat lifts and critical lengths}\label{app:flat}
Take square horizontal lattices for \(m=2,4\) and unit-edge hexagonal ones for \(m=3,6\). Let periodicity of a spinor under the three translations be \((-1)^i,(-1)^j,(-1)^k\). A right screw of angle \(2\pi/m\) and axial displacement \(L/m\) lifts as \(s\exp(-i\pi\sigma_3/m)\), \(s=\pm1\); thus \((-1)^k=-s^m\). The horizontal character must be rotation-invariant. There are four options for order two, two for order four, and one each for orders three and six. On the order-four manifold the two nonperiodic-horizontal lifts have paired spectra; on the Hantzsche--Wendt manifold all four spins have half-integral momenta in each direction and every orbit is paired. Only the fully periodic three-torus has a two-dimensional unshifted kernel.

For horizontal-periodic spin solve \(e^{2\pi iz/m}=s e^{\mp i\pi/m}\) along each axial eigenspinor. Both lines give the same axial tower with
\(c_*=(-\tfrac12+\tfrac m2\,\mathbf1_{s=-1})\bmod m\). The two Hurwitz zeta values at zero give the eta formula in Section~\ref{sec:catalog}. The pairs in \(s=+,-\) order are \((-1,1)\), \((-4/3,2/3)\), \((-3/2,1/2)\), \((-5/3,1/3)\) for \(m=2,3,4,6\). Pf\"affle's axial tower uses the opposite operator sign \cite[Thms.~5.4, 5.6]{pfaffle}; changing it once matches all eight nonzero values. The 20 remaining lifts have zero eta.

For a nonaxial mode the screw rotates the horizontal covector through its full orbit and leaves each sign once per orbit. Let \(z_0=k/2\) with translation bit \(k=0,1\). Comparing its minimum against the two axial levels gives splitting precisely when
\[
 L^2>(\min(c_*,m-c_*)^2-z_0^2)/\rho^2
\]
for the asymmetric lifts. At equality one should add multiplicities of any tied shells on each side. For \(m=6,s=-1\), \(\rho=2/\sqrt3,\ z_0=1/2,\ c_*=5/2\); the critical \(L^2\) is \(9/2\). Translation and screw lifts thus supply both the first-gap equality and its breakdown in a single metric family. Numerical charge units use \(q_\pm=4\vol|\lambda_\pm|\), including the holonomy quotient in the volume.
